\documentclass[12pt,a4paper]{article}

% ── Packages ──────────────────────────────────────────────
\usepackage[top=1.8cm, bottom=2cm, left=2cm, right=2cm]{geometry}
\usepackage{times}
\usepackage{enumitem}
\usepackage{tabularx}
\usepackage{graphicx}
\usepackage{multicol}
\usepackage{titlesec}
\usepackage{fancyhdr}
\usepackage{lastpage}
\usepackage{setspace}
\usepackage{ulem}       % for underline
\usepackage{soul}        % for highlighting
\usepackage{xcolor}
\usepackage{amssymb}

% ── No paragraph indent, small paragraph skip ────────────
\setlength{\parindent}{0pt}
\setlength{\parskip}{4pt}

% ── Disable default page numbering (we handle it manually) ─
\pagestyle{empty}

% ── Custom list styles ───────────────────────────────────
\newlist{romanenum}{enumerate}{1}
\setlist[romanenum]{label=\textbf{(\roman*)},leftmargin=1.2cm,itemsep=3pt,topsep=4pt}

\newlist{abcenum}{enumerate}{1}
\setlist[abcenum]{label=\textbf{(\alph*)},leftmargin=1.5cm,itemsep=1pt,topsep=2pt}

% ── Marks in right margin ────────────────────────────────
\newcommand{\qmarks}[1]{\hfill\textbf{[#1]}}

% ══════════════════════════════════════════════════════════
\begin{document}

% ── TITLE BLOCK ──────────────────────────────────────────
\begin{center}
    {\LARGE\bfseries ICSE BOARD PAPER - 2025}\\[4pt]
    {\Large\bfseries ENGLISH PAPER-2}\\[4pt]
    {\Large\bfseries [LITERATURE]}\\[4pt]
    {\Large\bfseries Class-10\textsuperscript{th}}\\[4pt]
    {\Large\bfseries (Solved)}
\end{center}

\vspace{6pt}

% ── Maximum Marks / Time line ────────────────────────────
\noindent\rule{\textwidth}{0.8pt}
\vspace{-6pt}

\noindent\begin{tabularx}{\textwidth}{@{}X r@{}}
\textbf{\textit{Maximum Marks: 80}} & \textbf{\textit{Time Allotted: Two Hours}}
\end{tabularx}

\vspace{-2pt}
\noindent\rule{\textwidth}{0.4pt}

\vspace{6pt}

% ── Instructions to Candidates ───────────────────────────
\textbf{Instructions to Candidates:}
\begin{enumerate}[leftmargin=0.8cm, itemsep=1pt, topsep=2pt]
    \item \textit{Answers to this Paper must be written on the paper provided separately.}
    \item \textit{You will \textbf{not} be allowed to write during the first \textbf{15} minutes.}
    \item \textit{This time is to be spent in reading the question paper.}
    \item \textit{\textbf{The time given at the head of this Paper is the time allowed for writing the answers.}}
\end{enumerate}

\vspace{2pt}
\noindent\rule{\textwidth}{0.4pt}
\vspace{4pt}

\begin{enumerate}[leftmargin=0.8cm, itemsep=1pt, topsep=2pt, start=5]
    \item \textit{The paper has \textbf{four} Sections.}
    \item \textit{\textbf{Section A} is compulsory -- \textbf{All} questions in \textbf{Section A} must be answered.}
    \item \textit{You must attempt \textbf{one} question from each of the \textbf{Sections B, C} and \textbf{D} and \textbf{one} other question from \textbf{any Section} of your choice.}
    \item \textit{The intended marks for questions or parts of questions are given in brackets~[~].}
\end{enumerate}

\vspace{2pt}
\noindent\rule{\textwidth}{0.4pt}

% ══════════════════════════════════════════════════════════
%                         SECTION A
% ══════════════════════════════════════════════════════════
\vspace{8pt}
\begin{center}
    {\large\bfseries SECTION A}\\[2pt]
    \textit{(Attempt \textbf{all} questions from this \textbf{Section}.)}
\end{center}

\vspace{4pt}

% ── Question 1 ───────────────────────────────────────────
\textbf{Question 1} \qmarks{16}

\textbf{Choose the correct answers to the questions from the given options.}\\
\textbf{\textit{(Do not copy the questions, write ONLY the correct answers.)}}

\begin{romanenum}

% (i)
\item What does Antony describe as, `\textit{thou bleeding piece of earth,}'?
\begin{abcenum}
    \item Caesar's robe that was stained with his blood
    \item The blood-soaked spot on which Caesar lies
    \item The bloodied pedestal on which Caesar fell
    \item Caesar's lifeless body covered in blood
\end{abcenum}

% (ii)
\item When Brutus says, `\textit{ambition's debt is paid}', he means that \rule{2cm}{0.4pt}.
\begin{abcenum}
    \item Caesar's greed for power and possession had resulted in his death
    \item Caesar had left seventy-five drachmas for every citizen of Rome
    \item Mark Antony would be permitted to speak at Caesar's funeral
    \item The conspirators would be punished for assassinating Caesar
\end{abcenum}

% (iii)
\item After Antony's speech following Caesar's assassination, the angry mob kills \rule{2cm}{0.4pt}.

\vspace{2pt}
\begin{tabularx}{\textwidth}{@{\hspace{1.5cm}}X X}
\textbf{(a)}~Cinna the conspirator & \textbf{(b)}~Cinna the poet \\
\textbf{(c)}~Decius Brutus & \textbf{(d)}~Cicero
\end{tabularx}

% (iv)
\item When Antony says, `\textit{He must be taught, and train'd, and bid go forth: A barren-spirited fellow;}' he refers to \rule{2cm}{0.4pt}.

\vspace{2pt}
\begin{tabularx}{\textwidth}{@{\hspace{1.5cm}}X X}
\textbf{(a)}~Octavius & \textbf{(b)}~Lucilius \\
\textbf{(c)}~Lucius & \textbf{(d)}~Lepidus
\end{tabularx}

% (v)
\item Who were the members of the Second Triumvirate, formed to rule over Rome after Caesar's assassination?

\vspace{2pt}
\begin{tabularx}{\textwidth}{@{\hspace{1.5cm}}X X}
\textbf{(a)}~Julius Caesar, Crassus, Pompey & \textbf{(b)}~Marcus Brutus, Caius Cassius, Casca \\
\textbf{(c)}~Mark Antony, Octavius Caesar, Lepidus & \textbf{(d)}~Octavius Caesar, Julius Caesar, Mark Antony
\end{tabularx}

% (vi)
\item Which of the following words best describes Cassius's mood when he says:
\begin{center}
\textit{`Come, Antony, and young Octavius, come,}\\
\textit{Revenge yourselves alone on Cassius,'}
\end{center}

\vspace{2pt}
\begin{tabularx}{\textwidth}{@{\hspace{1.5cm}}X X}
\textbf{(a)}~anxious & \textbf{(b)}~regretful \\
\textbf{(c)}~fearful & \textbf{(d)}~frustrated
\end{tabularx}

% (vii)
\item In the poem, `\textit{Haunted Houses}', what does the speaker see at his fireside that the stranger cannot?
\begin{abcenum}
    \item the speaker only sees what is physically present
    \item the stranger can see the spirits while the speaker cannot
    \item the speaker sees both the present and the past
    \item the stranger sees both the present and the future
\end{abcenum}

% (viii)
\item In the poem, `\textit{The Glove and the Lions}', who does Count de Lorge `\textit{sigh for}'?
\begin{abcenum}
    \item The lions that ramped and roared
    \item King Francis who loved a royal sport
    \item The nobles who filled the benches
    \item The beauteous lively dame
\end{abcenum}

% (ix)
\item Select the option that shows the correct relationship between Statements (1) and (2) from Maya Angelou's poem, `\textit{When Great Trees Fall}':

\textbf{Statement 1:} When great trees fall, the world is left permanently unstable.

\textbf{Statement 2:} In the poem, the natural world is depicted as reacting with fear and uncertainty when great trees fall.

\vspace{2pt}
\begin{tabularx}{\textwidth}{@{\hspace{1.5cm}}X X}
\textbf{(a)}~(1) is false, but (2) is true. & \textbf{(b)}~Both (1) and (2) are false. \\
\textbf{(c)}~(1) is true, but (2) is false. & \textbf{(d)}~Both (1) and (2) are true.
\end{tabularx}

% (x)
\item The poem, `\textit{A Considerable Speck}', expresses Frost's \rule{3cm}{0.4pt}.

\vspace{2pt}
\begin{tabularx}{\textwidth}{@{\hspace{1.5cm}}X X}
\textbf{(a)}~love for the little organism & \textbf{(b)}~respect for intelligent life \\
\textbf{(c)}~indifference to the tiny creature & \textbf{(d)}~anger at the mite's escape
\end{tabularx}

% (xi)
\item Which of the given options contains the figure of speech that appears in the following line from the poem, `\textit{The Power of Music}'?

\begin{center}
\textit{And in the sky the feathered fly turn turtle while}\\
\textit{They're winging,}
\end{center}

\begin{abcenum}
    \item You're one month on in the middle of May
    \item Rainbow-tinted circles of light
    \item And life is too much like a pathless wood
    \item The wind lies asleep in the arms of the dawn
\end{abcenum}

% (xii)
\item In the short story, `\textit{With the Photographer}', the narrator asks the photographer, ``\textit{Is it me?}'' because \rule{2cm}{0.4pt}.
\begin{abcenum}
    \item He is delighted with the photographer's skill
    \item He looks very handsome in the photograph
    \item He is unable to recognise his own face
    \item He is ashamed of how he looks in the picture
\end{abcenum}

% (xiii)
\item In the short story, `\textit{The Elevator}', which of the following does Martin's father \textbf{\textit{NOT}} do when Martin expresses his fear of the elevator?

\vspace{2pt}
\begin{tabularx}{\textwidth}{@{\hspace{1.5cm}}X X}
\textbf{(a)}~He dismisses Martin's concerns & \textbf{(b)}~He encourages Martin to use the stairs \\
\textbf{(c)}~He tells Martin to grow up and be brave & \textbf{(d)}~He watches TV and ignores Martin's fear
\end{tabularx}

% (xiv)
\item Choose the option that lists the sequence of events from Ray Bradbury's short story, `\textit{The Pedestrian}':
\begin{enumerate}[leftmargin=0.8cm, itemsep=1pt, label=\textbf{\arabic*.}]
    \item One night, as he was nearing home, he was stopped by a police car.
    \item Since his answers were considered odd and unacceptable, he was forced to enter the car and taken away.
    \item Leonard Mead loved to walk through the empty streets at night.
    \item A metallic voice from the car asked him a series of questions.
\end{enumerate}

\vspace{2pt}
\begin{tabularx}{\textwidth}{@{\hspace{1.5cm}}X X}
\textbf{(a)}~2, 1, 3, 4 & \textbf{(b)}~1, 3, 4, 2 \\
\textbf{(c)}~3, 1, 4, 2 & \textbf{(d)}~4, 1, 2, 3
\end{tabularx}

% (xv)
\item Where was Adjoa born?

\vspace{2pt}
\begin{tabularx}{\textwidth}{@{\hspace{1.5cm}}X X}
\textbf{(a)}~Nigeria & \textbf{(b)}~Ghana \\
\textbf{(c)}~Ethiopia & \textbf{(d)}~Kenya
\end{tabularx}

% (xvi)
\item Why did M.\ Hamel have to leave the school after forty years of service?
\begin{abcenum}
    \item He was \textit{not} allowed to teach French any longer
    \item He had grown tired of his job as a school teacher
    \item He was thought to be too strict in his ways
    \item He wanted to retire and take up farming
\end{abcenum}

\end{romanenum}

\newpage

% ══════════════════════════════════════════════════════════
%                         SECTION B
% ══════════════════════════════════════════════════════════
\begin{center}
    {\large\bfseries SECTION B}\hfill\textbf{[10]}\\[2pt]
    \textit{(Answer \textbf{one or more} questions from this \textbf{Section}.)}
\end{center}

\vspace{4pt}
\begin{center}
    {\large\bfseries DRAMA}\\[2pt]
    \textbf{(Julius Caesar by William Shakespeare)}
\end{center}

\vspace{6pt}

% ── Question 2 ───────────────────────────────────────────
\textbf{Question 2}\\
Read the extract from `\textit{Julius Caesar}' Act 3, Scene 2, given below and answer the questions that follow:

\vspace{4pt}
\begin{tabular}{@{}l@{\hspace{1cm}}l@{}}
\textbf{Antony:} & \textit{Friends, Romans, countrymen, lend me your ears;} \\
 & \textit{I come to bury Caesar, not to praise him.} \\
 & \textit{The evil that men do lives after them,} \\
 & \textit{The good is oft interred with their bones;} \\
 & \textit{So let it be with Caesar. The noble Brutus} \\
 & \textit{Hath told you Caesar was ambitious;} \\
 & \textit{If it were so, it was a grievous fault,} \\
 & \textit{And grievously hath Caesar answer'd it.} \\
 & \textit{Here, under leave of Brutus and the rest--} \\
 & \textit{For Brutus is an honourable man;} \\
 & \textit{So are they all, all honourable men--} \\
 & \textit{Come I to speak in Caesar's funeral.}
\end{tabular}

\vspace{8pt}

\begin{romanenum}
    \item What does Antony say he is there for?\\
          What does he say he is \textbf{\textit{not}} there for?\\
          What do you think he is actually there for? \qmarks{3}

    \item What were the \textbf{\textit{three}} conditions that Brutus had laid down before allowing Antony to speak to the citizens of Rome? \qmarks{3}

    \item List the \textbf{\textit{three}} arguments that Antony uses immediately after he speaks these lines to prove conclusively that Julius Caesar was not ambitious. \qmarks{3}

    \item Antony repeatedly uses certain words in his speech to describe Brutus. What are they? Why does he do this? \qmarks{3}

    \item Which \textbf{\textit{one}} argument of Antony's do you think had the greatest impact on his listeners? Give a reason to support your answer. \qmarks{4}

    What were the citizens' feelings towards Antony before he began his speech?\\
    How do their feelings towards him change at the end of his speech?
\end{romanenum}

\vspace{10pt}

% ── Question 3 ───────────────────────────────────────────
\textbf{Question 3}\\
Read the extract from `\textit{Julius Caesar}' Act 5, Scene 1, given below and answer the questions that follow:

\vspace{4pt}
\begin{tabular}{@{}l@{\hspace{1cm}}l@{}}
\textbf{Cassius:} & \textit{Then, if we lose this battle,} \\
 & \textit{You are contented to be led in triumph} \\
 & \textit{Through the streets of Rome?}
\end{tabular}

\vspace{4pt}
\begin{tabular}{@{}l@{\hspace{1cm}}l@{}}
\textbf{Brutus:} & \textit{No, Cassius, no: think not, thou noble Roman,} \\
 & \textit{That ever Brutus will go bound to Rome;} \\
 & \textit{He bears too great a mind. But this same day} \\
 & \textit{Must end that work the ides of March begun;} \\
 & \textit{And whether we shall meet again I know not.} \\
 & \textit{Therefore our everlasting farewell take:} \\
 & \textit{Forever, and forever, farewell, Cassius.} \\
 & \textit{If we do meet again, why, we shall smile;} \\
 & \textit{If not, why then, this parting was well made.}
\end{tabular}

\vspace{8pt}

\begin{romanenum}
    \item Why does Brutus say, `\textit{No, Cassius, no}'? \qmarks{3}

    What conditions may force Brutus to go bound to Rome?

    Brutus says, `\textit{And whether we shall meet again I know not.}' What do these words imply?

    \item To what does Brutus refer when he says, `\textit{the work which the ides of March begun}'? \qmarks{3}

    How was `\textit{that work}' begun?

    What political change would take place in Rome if Brutus and Cassius lose this battle?

    \item Earlier in this scene, Cassius had confided to Messala a strange occurrence that he had observed when his army was on its way from Sardis to Philippi. \qmarks{3}

    Describe this strange occurrence.

    \item How does Brutus die? \qmarks{3}

    How does his manner of dying contradict the philosophy by which he had lived his life?

    \item At the end of the play, Antony calls Brutus, `\textit{the noblest Roman of them all}'. Why does he say this? \qmarks{3}

    What does this reveal of Antony's character?
\end{romanenum}

\newpage

% ══════════════════════════════════════════════════════════
%                         SECTION C
% ══════════════════════════════════════════════════════════
\begin{center}
    {\large\bfseries SECTION C}\\[2pt]
    \textit{(Answer \textbf{one or more} questions from this \textbf{Section}.)}
\end{center}

\vspace{4pt}
\begin{center}
    {\large\bfseries PROSE - SHORT STORIES}\\[2pt]
    \textbf{(Treasure Trove -- A Collection of Poems and Short Stories)}
\end{center}

\vspace{6pt}

% ── Question 4 ───────────────────────────────────────────
\textbf{Question 4}\\
Read the following extract from William Sleater's short story, `\textbf{\textit{The Elevator}}' and answer the questions that follow:

\vspace{4pt}
\begin{quote}
\textit{Martin felt nervous when he got back to the building after school. But why should he be afraid of an old lady? He felt ashamed of himself. He pressed the button and stepped into the elevator, hoping that it would not stop, but it stopped on the third floor. Martin watched the door slide open \ldots}
\end{quote}

\begin{romanenum}
    \item Who entered the elevator when the door slid open? Describe this person. \qmarks{3}

    \item Why was Martin afraid of using this elevator? \qmarks{3}

    \item What led to Martin's fall down the stairs? \qmarks{3}

    What did it result in?

    \item How does the story end? \qmarks{3}

    Bring out the element of horror in the ending.

    \item What was Martin's father's opinion of him? \qmarks{4}

    What does this reveal to us about his father's character?
\end{romanenum}

\vspace{10pt}

% ── Question 5 ───────────────────────────────────────────
\textbf{Question 5}\\
Read the following extract from Alphonse Daudet's short story, `\textbf{\textit{The Last Lesson}}' and answer the questions that follow:

\vspace{4pt}
\begin{quote}
\textit{``My children, this is the last lesson I shall give you. The order has come from Berlin to teach only German in the schools of Alsace and Lorraine. The new master comes tomorrow. This is your last French lesson. I want you to be very attentive.''}

\textit{What a thunder-clap these words to me!}

\textit{Oh the wretches; \textbf{that} was what they had put up at the town-hall!}
\end{quote}

\begin{romanenum}
    \item What does the word, `\textbf{\textit{that}}' in the extract refer to? \qmarks{3}

    What other `bad news' had the villagers received from the bulletin board outside the town-hall?

    \item How does the narrator describe the daily bustle at the start of a typical school day? \qmarks{3}

    \item What was different on that morning when he arrived late for school? \qmarks{3}

    \item What thoughts filled the narrator's head when he heard the above announcement? \qmarks{3}

    \item What does M.\ Hamel urge his listeners never to forget? \qmarks{4}

    M.\ Hamel writes a few words on the board at the end of the lesson. What were they?

    What does he hope to inspire through his words and actions on that day?
\end{romanenum}

\newpage

% ══════════════════════════════════════════════════════════
%                         SECTION D
% ══════════════════════════════════════════════════════════
\begin{center}
    {\large\bfseries SECTION D}\\[2pt]
    \textit{(Answer \textbf{one or more} questions from this \textbf{Section}.)}
\end{center}

\vspace{4pt}
\begin{center}
    {\large\bfseries POETRY}\\[2pt]
    \textbf{(Treasure Trove -- A Collection of Poems and Short Stories)}
\end{center}

\vspace{6pt}

% ── Question 6 ───────────────────────────────────────────
\textbf{Question 6}\\
Read the following extract from Leigh Hunt's poem, `\textbf{\textit{The Glove and the Lions}}' and answer the questions that follow:

\vspace{4pt}
\begin{center}
\textit{And truly 'twas a gallant thing to see that crowning show,}\\
\textit{Valour and love, and a king above, and the royal beasts below.}
\end{center}

\begin{romanenum}
    \item Describe the scene at the beginning of the poem. \qmarks{3}

    \item ``Leigh Hunt uses vivid sound and visual images to describe the contest between the royal beasts in the pit.'' Justify with close reference to the text. \qmarks{3}

    \item Who was the `\textit{beauteous lively dame}' mentioned in the poem? \qmarks{3}

    What did she do?

    What prompted her to do this?

    \item Explain the following lines in your own words: \qmarks{3}
    \begin{abcenum}
        \item \textit{`smiling lips and sharp bright eyes which always seemed the same'}
        \item \textit{`the occasion is divine'}
        \item \textit{`Faith, gentlemen, we're better here than there'}
    \end{abcenum}

    \item When King Francis exclaims, ``\textit{rightly done!}'', what action of Count de Lorge does he applaud? \qmarks{4}

    Why do you think the Count behaved the way he did?
\end{romanenum}

\vspace{10pt}

% ── Question 7 ───────────────────────────────────────────
\textbf{Question 7}\\
Read the following extract from Robert Frost's poem, `\textbf{\textit{A Considerable Speck}}' and answer the questions that follow:

\vspace{4pt}
\begin{center}
\textit{This was no dust speck by my breathing blown,}\\
\textit{But unmistakably a living mite}\\
\textit{With inclinations it could call its own.}
\end{center}

\begin{romanenum}
    \item What did the narrator first imagine the `speck' to be? \qmarks{3}

    What does his use of the words `speck' and `mite' tell us about it?

    \item What was the narrator doing when he spotted the speck? \qmarks{3}

    Which \textbf{two} words from the poem helped you come to this conclusion?

    \item What made him realise that he was dealing with `an intelligence'? \qmarks{3}

    \item Why did the narrator conclude that the mite had feet? \qmarks{3}

    \item Describe the narrator's initial response to the speck. \qmarks{4}

    What did he eventually decide to do about it?

    Why does he make this decision?

    What does this decision reveal of his character?
\end{romanenum}

\vspace{20pt}
\begin{center}
$\blacksquare$~$\blacksquare$
\end{center}

\end{document}
